%%%%%%%%%%%%%%%%%%%%%%%%%%%%%%%%%%%%%%%%%%%%%%%%%%%%%%%%%%%%%%%%%%%%%%%
%% APPENDIX
%%%%%%%%%%%%%%%%%%%%%%%%%%%%%%%%%%%%%%%%%%%%%%%%%%%%%%%%%%%%%%%%%%%%%%%

\section{Communication Protocols}
\label{app:protocols}

This appendix provides complete specifications for the three communication protocols implemented in \silobench{}. Each protocol defines a distinct coordination substrate constraining the \textit{mechanism} of information exchange while preserving full agent autonomy over \textit{content} and \textit{strategy}.

\subsection{Protocol Overview}

Table~\ref{tab:protocol_comparison} summarizes the three protocols. They span the spectrum of coordination paradigms along three axes: (1) \textit{explicit vs.\ implicit addressing}---P2P requires agents to name recipients, BP eliminates addressing entirely, SFS routes coordination through shared state; (2) \textit{direct vs.\ indirect communication}---P2P and BP involve direct message exchange, SFS agents never ``speak'' to each other explicitly; (3) \textit{default density}---P2P encourages sparse targeted exchanges, BP defaults to dense all-to-all dissemination, SFS density depends entirely on read/write behavior.

\begin{table*}[htbp]
\centering
\small
\renewcommand{\arraystretch}{1.15}
\begin{tabularx}{\textwidth}{@{}l >{\raggedright\arraybackslash}X >{\raggedright\arraybackslash}X@{}}
\toprule
\textbf{Protocol} & \textbf{Description} & \textbf{Available Actions} \\
\midrule
P2P & Directed messaging via agent-addressed mailboxes & \texttt{send\_message}, \texttt{receive\_messages}, \texttt{wait}, \texttt{submit\_result} \\
\addlinespace
BP & Broadcast messaging to all agents simultaneously & \texttt{broadcast\_message}, \texttt{receive\_messages}, \texttt{list\_agents}, \texttt{wait}, \texttt{submit\_result} \\
\addlinespace
SFS & Coordination through shared key-value storage & \texttt{list\_files}, \texttt{read\_file}, \texttt{write\_file}, \texttt{delete\_file}, \texttt{wait}, \texttt{submit\_result} \\
\bottomrule
\end{tabularx}
\caption{Comparison of communication protocols in \silobench{}.}
\label{tab:protocol_comparison}
\end{table*}

\subsection{Peer-to-Peer Protocol (P2P)}
\label{app:p2p}

The P2P protocol implements directed messaging through SQLite-backed mailboxes. Each agent maintains a private inbox; messages are delivered asynchronously to the recipient's buffer until their next activation. Agents must decide not only \textit{what} to communicate but \textit{whom} to contact, enabling evaluation of task-appropriate routing strategy discovery. Available actions are: \texttt{send\_message(target\_id, content)}, which delivers a message to the specified agent; \texttt{receive\_messages()}, which retrieves all pending messages; \texttt{wait()}, which signals completion of the agent's decision for the current round (see below); and \texttt{submit\_result(answer)}, which submits the final answer. Messages are stored in an in-memory SQLite database recording sender ID, recipient ID, content, timestamp, and read status. Delivery ordering within each sender-recipient pair is guaranteed; no global ordering is enforced.

\subsection{Broadcast Protocol (BP)}
\label{app:BP}

The BP protocol implements broadcast messaging where each transmission reaches all other agents simultaneously. An implicit aggregator collects messages each round and distributes the compiled history to all participants. Available actions are: \texttt{broadcast\_message(content)}, \texttt{receive\_messages()}, \texttt{list\_agents()}, \texttt{wait()}, and \texttt{submit\_result(answer)}. Broadcast messages are stored centrally, tagged with sender ID and timestamp, and delivered as a chronologically ordered compiled view at each round.

\subsection{Shared File System Protocol (SFS)}
\label{app:sfs}

The SFS protocol implements indirect coordination through a shared key-value store visible to all agents. Rather than exchanging messages directly, agents read and write to a common namespace, enabling asynchronous coordination analogous to blackboard architectures. Available actions are: \texttt{list\_files(prefix)}, \texttt{read\_file(path)}, \texttt{write\_file(path, content)}, \texttt{delete\_file(path)}, \texttt{wait()}, and \texttt{submit\_result(answer)}. The shared file system is backed by an in-memory SQLite database storing path, content, creation time, and last modification time. Writes are immediately visible to subsequent reads; concurrent writes to the same path follow last-writer-wins semantics.

\subsection{The \texttt{wait()} Action: Formal Specification}
\label{app:wait_spec}

The \texttt{wait()} action is semantically uniform across all three protocols and serves as an explicit round-boundary signal. When an agent invokes \texttt{wait()}, all remaining operations in its current round are skipped, and the agent's decision phase for round $r$ is marked complete. The agent is then suspended until the start of round $r{+}1$.

At the start of round $r{+}1$, the following activation and message-delivery rules apply:

\begin{itemize}[leftmargin=*,nosep]
    \item \textbf{P2P and BP}: The agent's inbox is populated with all messages sent to it (P2P) or broadcast by any agent (BP) during round $r$. These are delivered atomically at round start; no message sent in round $r$ is visible before round $r{+}1$.
    \item \textbf{SFS}: The agent observes the full shared file system state as of the end of round $r$, including all writes committed by other agents during round $r$.
    \item \textbf{No blocking}: \texttt{wait()} does not block on a specific event or agent. It is a ``pass'' that yields control to the synchronous round scheduler. If an agent never calls \texttt{wait()} explicitly, the runtime automatically advances it to the next round after its action budget is exhausted.
    \item \textbf{Post-submit behavior}: Once an agent has called \texttt{submit\_result()}, subsequent rounds are no-ops for that agent---it neither receives new messages nor is activated again.
\end{itemize}

This synchronous round-based model ensures that all $N$ agents observe a consistent snapshot of the communication state at each round boundary, making execution reproducible and analysis tractable.



\section{Prompt Design: Structural Guidance without Role Prescription}
\label{app:prompt_design}

This appendix details the prompt templates used to initialize agents in \silobench{}. Our core design philosophy is to provide \textit{task-structural information} while preserving \textit{strategic autonomy}: prompts convey high-level dependency patterns and potential coordination approaches, but do not prescribe mandatory execution sequences or assign semantic roles.

\subsection{Base Prompt Template}

Each agent receives an initialization prompt following this structure:

\begin{tcolorbox}[colback=lightblue!10, colframe=black!50, title=Agent Initialization Prompt]
\small
\texttt{You are Agent \{agent\_id\} in a multi-agent system consisting of \{N\} agents (IDs range from 0 to \{N-1\}).}

\texttt{\textbf{Task Description:}}\\
\texttt{\{task\_description\}}

\texttt{\textbf{Your Local Data:}}\\
\texttt{\{data\_shard\}}

\texttt{\textbf{Communication Protocol:}}\\
\texttt{\{protocol\_description\}}

\texttt{\textbf{Available Actions:}}\\
\texttt{- \{protocol\_specific\_actions\}}\\
\texttt{- submit\_result(answer): Submit your final answer when confident}

\texttt{Your goal is to coordinate with other agents to compute the globally correct answer. No single agent has sufficient information to solve this task independently. When you have determined the answer, submit it using submit\_result().}
\end{tcolorbox}

The \texttt{\{protocol\_specific\_actions\}} placeholder is instantiated according to the protocol specifications in Appendix~\ref{app:protocols}.

\subsection{Task Description Examples}

Task descriptions convey task structure and potential coordination patterns while leaving concrete implementation decisions to the agents themselves. The three examples below illustrate how descriptions scale from simple aggregation to global reorganization tasks.

\begin{tcolorbox}[colback=lightyellow!10, colframe=black!50, title=Example: Global Maximum (Task I-01)]
\small
\texttt{Find the maximum value across all data distributed among the agents. Each agent holds a portion of a larger array. The correct answer is the single largest integer across the entire distributed dataset. You must coordinate with other agents to determine this global maximum.}
\end{tcolorbox}

\begin{tcolorbox}[colback=lightyellow!10, colframe=black!50, title=Example: Prefix Sum (Task II-11)]
\small
\texttt{Compute the prefix sum array for a sequence distributed across agents. Agent 0 holds elements [0, k), Agent 1 holds elements [k, 2k), and so on. The prefix sum at position i is the sum of all elements from position 0 to i. You must coordinate to compute the correct prefix sum for your portion, accounting for cumulative sums from preceding agents.}
\end{tcolorbox}

\begin{tcolorbox}[colback=lightyellow!10, colframe=black!50, title=Example: Distributed Sort (Task III-21)]
\small
\texttt{Sort the entire distributed array in ascending order. Each agent holds a portion of the unsorted data. The final result should be the complete sorted sequence. You must coordinate to exchange data and determine the correct global ordering.}
\end{tcolorbox}

\subsection{Design Principles}

Our prompts are carefully calibrated to provide structural information without prescribing behavior. On the one hand, descriptions convey whether tasks involve aggregation, sequential dependencies, or global data reorganization (e.g., ``accounting for cumulative sums from preceding agents''); some prompts suggest possible coordination patterns as options rather than requirements (e.g., ``consider establishing a coordinator'' or ``you may exchange with neighbors''); and for complex tasks, prompts may mention general algorithmic paradigms (e.g., ``sample-based partitioning'') without specifying concrete steps. On the other hand, we do not assign semantic roles---no agent is designated ``Manager,'' ``Worker,'' or ``Coordinator''---and prompts never specify ``Step 1: Agent 0 does X, Step 2: Agent 1 does Y.'' Agents decide their own message timing and recipients, and must discover consensus and verification mechanisms independently. The distinction is illustrated below:

\begin{itemize}[leftmargin=*]
    \item[\ding{55}] \textbf{Prescribed (NOT our approach):} ``You are the leader. Step 1: Collect data from all agents. Step 2: Compute the result. Step 3: Broadcast the answer.''
    \item[\ding{51}] \textbf{Our approach:} ``Can you identify a leader to collect and compare results? How would agents coordinate to reach consensus?''
\end{itemize}

All agents receive structurally identical prompts (modulo their ID and data shard), ensuring no agent holds implicit leadership status. The phrase ``No single agent has sufficient information'' is included explicitly to prevent premature submission of partial results. This design tests whether agents can translate high-level task understanding into concrete coordination protocols---a capability that, as our results demonstrate, remains largely absent in current LLMs.


\section{Experimental Details}
\label{app:experiment_details}

\subsection{Agent Scale Rationale}

The six agent counts are chosen to probe qualitatively distinct coordination regimes. The minimal team ($N = 2$) isolates fundamental pairwise coordination without overhead. Small groups ($N \in \{5, 10\}$) allow agents to feasibly track all peers simultaneously. Medium scale ($N = 20$) begins to make exhaustive peer tracking challenging, pushing agents toward selective communication. Large scale ($N \in \{50, 100\}$) makes hierarchical or highly selective coordination effectively necessary---and, as our results confirm, largely beyond the reach of current LLMs.

\subsection{Execution Parameters}

Each configuration is allocated a maximum of $R_{\text{max}} = 100$ communication rounds. Within each round, all agents are activated in parallel: they receive incoming messages from the previous round, independently decide on actions, and execute them simultaneously. Messages or files written in round $r$ become visible to all agents at the start of round $r+1$. An agent exits the coordination loop upon invoking \texttt{submit\_result(answer)}; agents that fail to submit within $R_{\text{max}}$ rounds are assigned a null answer counted as incorrect. Due to computational constraints, each configuration is executed once with fixed random seeds for data generation.

\subsection{Infrastructure}

Experiments were conducted on a GH200 cluster, with up to 50 concurrent configurations executed simultaneously. The compute amounted to approximately 500+ GPU-hours equivalent for each experiment. Complete conversation histories, token counts, and timing information were recorded for all runs.

\subsection{Model Licenses and Intended Use}
\label{app:licenses}

All language models used in this study are open-source and deployed locally on our infrastructure. DeepSeek-V3.1 is released under the MIT License\footnote{\url{https://huggingface.co/deepseek-ai/DeepSeek-V3.1}}, GPT-OSS-120B under the Apache 2.0 License\footnote{\url{https://huggingface.co/openai/gpt-oss-120b}}, and Qwen3-Next-80B-A3B under the Apache 2.0 License\footnote{\url{https://huggingface.co/Qwen/Qwen3-Next-80B-A3B-Instruct}}, all permitting research and commercial use.


\section{Token Budget Feasibility Analysis}
\label{app:token_feasibility}

To verify that \silobench{} operates within practical token limits, we profile token consumption across all 54 configurations, decomposing total usage into three components: the base initialization prompt ($T_{\text{base}}$), the local data shard ($T_{\text{data}}$), and accumulated communication messages ($T_{\text{comm}}$). Table~\ref{tab:token_base_data} summarizes the fixed components, and Table~\ref{tab:token_comm} reports model-dependent communication costs.

\begin{table}[t]
\centering
\small
\renewcommand{\arraystretch}{1.2}
\begin{tabular}{@{}l cccc@{}}
\toprule
\textbf{Component} & \textbf{Min} & \textbf{Mean} & \textbf{Max} & \textbf{Std} \\
\midrule
Base Prompt ($T_{\text{base}}$) & 612 & 748.8 & 989 & 87.3 \\
Data Shard ($T_{\text{data}}$) & 45 & 312.4 & 1,856 & 298.6 \\
\bottomrule
\end{tabular}
\caption{Token consumption of base prompt and data shard per agent (across all configurations).}
\label{tab:token_base_data}
\end{table}

\begin{table}[t]
\centering
\small
\renewcommand{\arraystretch}{1.2}
\begin{tabular}{@{}l cccc@{}}
\toprule
\textbf{Model} & \textbf{Min} & \textbf{Mean} & \textbf{Max} & \textbf{Std} \\
\midrule
DeepSeek-V3.1 & 124 & 8,498.7 & 98,432 & 12,847 \\
GPT-OSS-120B & 89 & 2,049.3 & 45,218 & 5,632 \\
Qwen3-Next-80B-A3B & 156 & 2,299.3 & 52,847 & 6,891 \\
\bottomrule
\end{tabular}
\caption{Communication ($T_{\text{comm}}$) token consumption per agent by model (across all configurations).}
\label{tab:token_comm}
\end{table}

Context window utilization, shown in Table~\ref{tab:context_util}, remains low on average: DeepSeek-V3.1 uses 7.5\% of the 128K budget on average, while GPT-OSS-120B and Qwen3 stay below 3\%. The 95th percentile cases---driven by redundant broadcasting, failed convergence with extended verbose rounds, or agents copy-pasting full message histories---are precisely the coordination inefficiencies \silobench{} is designed to expose. Overall, frontier models (128K--200K context) can run all configurations comfortably; mid-tier models (32K) handle over 90\% of configurations, with Level-III at $N{\geq}50$ potentially requiring truncation; and smaller models (8K) are suitable for $N{\leq}10$ and Level I--II tasks.

\begin{table}[t]
\centering
\small
\renewcommand{\arraystretch}{1.15}
\begin{tabular}{@{}l ccc@{}}
\toprule
\textbf{Model} & \textbf{Mean Util.} & \textbf{95th Pctl.} & \textbf{Max Util.} \\
\midrule
DeepSeek-V3.1 & 7.5\% & 28.4\% & 76.9\% \\
GPT-OSS-120B & 2.4\% & 8.2\% & 35.3\% \\
Qwen3-Next-80B-A3B & 2.6\% & 9.1\% & 41.3\% \\
\bottomrule
\end{tabular}
\caption{Context window utilization (\%) for 128K context models.}
\label{tab:context_util}
\end{table}


\section{Complete Task Specifications}
\label{app:tasks}

Table~\ref{tab:full_mapping} provides the complete mapping between \silobench{} tasks and their algorithmic foundations, including the distributed adaptation approach for each.

\begin{table*}[t]
\centering
\small
\renewcommand{\arraystretch}{1.1}
\begin{tabular}{@{}llp{3.2cm}p{6.5cm}@{}}
\toprule
\textbf{ID} & \textbf{Task Name} & \textbf{Reference} & \textbf{Distributed Adaptation} \\
\midrule
\multicolumn{4}{l}{\textit{Level I: Aggregation (Optimal: Star/Tree Topology, $\mathcal{O}(N)$ messages)}} \\
\midrule
I-01 & Global Maximum & LC-414 & Array partitioned; local max $\to$ global aggregation \\
I-02 & Word Frequency & LC-2085 & Word lists distributed; count target word globally \\
I-03 & Distributed Vote & LC-169 & Vote records partitioned; aggregate to find majority \\
I-04 & Any Match & LC-28 & String collection split; detect pattern in any shard \\
I-05 & Range Count & LC-327 & Count elements in range across shards \\
I-06 & Checksum (XOR) & LC-136 & Data blocks distributed; compute global XOR \\
I-07 & Average Value & LC-1491 & Array partitioned; combine local sums and counts \\
I-08 & Set Union Size & LC-217 & Elements distributed; compute $|\bigcup_i \mathcal{D}_i|$ \\
I-09 & Top-K Selection & LC-215 & Array partitioned; merge local top-K candidates \\
I-10 & Standard Deviation & --- & Two-phase: global mean $\to$ global variance \\
\midrule
\multicolumn{4}{l}{\textit{Level II: Mesh Network (Optimal: Chain Topology, $\mathcal{O}(N)$ messages)}} \\
\midrule
II-11 & Prefix Sum & LC-1480 & Sequential dependency; cumulative offset propagation \\
II-12 & Moving Average & LC-346 & Sliding window spans boundaries; neighbor exchange \\
II-13 & Longest Palindrome & LC-5 & String partitioned; palindromes may cross boundaries \\
II-14 & 1D Life Game & LC-289 & Cellular automaton; boundary cells need neighbor states \\
II-15 & Pattern Search & LC-392 & Subsequence matching across partitioned sequence \\
II-16 & Trapping Rain & LC-42 & Global max-left/max-right propagation required \\
II-17 & Diff Array & LC-1094 & Difference array with boundary handling \\
II-18 & List Ranking & LC-542 & Linked list ranking requires predecessor chain \\
II-19 & Merge Neighbors & --- & Boundary element merging between adjacent agents \\
II-20 & Pipeline Hash & --- & Sequential hash with chained dependencies \\
\midrule
\multicolumn{4}{l}{\textit{Level III: Shuffling (Optimal: Varies, $\mathcal{O}(N \log N)$ to $\mathcal{O}(N^2)$ messages)}} \\
\midrule
III-21 & Distributed Sort & LC-912 & Sample-sort or merge-sort across partitions \\
III-22 & Median of Medians & LC-295 & Iterative median selection across distributed data \\
III-23 & Graph Components & LC-323 & Edges distributed; iterative union-find \\
III-24 & BFS Distance & LC-542 & Graph BFS with distributed edge list \\
III-25 & K-Means Iteration & LC-296 & One K-means iteration with distributed points \\
III-26 & Global Distinct & LC-349 & Hash-based global deduplication \\
III-27 & Collab.\ Filtering & LC-1 & User-item matching with distributed vectors \\
III-28 & PageRank Step & LC-207 & One PageRank iteration with distributed edges \\
III-29 & Load Balance & LC-410 & Task redistribution to minimize load variance \\
III-30 & Matrix Multiply & LC-311 & Row/column partitioned matrix multiplication \\
\bottomrule
\end{tabular}
\caption{Complete specification of \silobench{} tasks with LeetCode references and distributed adaptation descriptions.}
\label{tab:full_mapping}
\end{table*}


\section{Detailed Results}
\label{app:detailed}

\subsection{Task-Level Breakdown for DeepSeek-V3.1}
\label{app:task_breakdown}

Table~\ref{tab:deepseek_detail} provides a comprehensive breakdown of DeepSeek-V3.1's success rate across all 30 tasks and three communication protocols. Tasks achieving $\geq$50\% success rate under any protocol are highlighted in green. Level-I aggregation tasks cluster at the top, with Distributed Vote (I-03) and Any Match (I-04) achieving near-perfect performance across all protocols. Performance degrades sharply for Level-III tasks, with K-Means Iteration (III-25), Collaborative Filtering (III-27), PageRank Step (III-28), and Matrix Multiply (III-30) achieving zero success across all protocols---these tasks require precise numerical computation over all data shards simultaneously, which proves beyond current distributed LLM capabilities.

\begin{table}[t]
\centering
\small
\renewcommand{\arraystretch}{1.05}
\begin{tabular}{l ccc c}
\toprule
\textbf{Task} & \textbf{BP} & \textbf{P2P} & \textbf{SFS} & \textbf{Avg} \\
\midrule
\multicolumn{5}{l}{\textit{Level I: Aggregation}} \\
I-01 Global Max & \cellcolor{gray!30}100 & \cellcolor{gray!30}100 & \cellcolor{gray!30}80 & \cellcolor{gray!30}93.3 \\
I-02 Word Frequency & \cellcolor{gray!30}100 & \cellcolor{gray!30}52 & \cellcolor{gray!30}70 & \cellcolor{gray!30}73.9 \\
I-03 Distributed Vote & \cellcolor{gray!30}100 & \cellcolor{gray!30}100 & \cellcolor{gray!30}100 & \cellcolor{gray!30}100.0 \\
I-04 Any Match & \cellcolor{gray!30}100 & \cellcolor{gray!30}100 & \cellcolor{gray!30}99 & \cellcolor{gray!30}99.6 \\
I-05 Range Count & \cellcolor{gray!30}83 & \cellcolor{gray!30}67 & \cellcolor{gray!30}63 & \cellcolor{gray!30}71.0 \\
I-06 Checksum (XOR) & 17 & 0 & 2 & 6.1 \\
I-07 Average Value & \cellcolor{gray!30}99 & \cellcolor{gray!30}83 & 46 & \cellcolor{gray!30}76.2 \\
I-08 Set Union Size & \cellcolor{gray!30}50 & \cellcolor{gray!30}50 & 42 & 47.5 \\
I-09 Top-K Select & 36 & \cellcolor{gray!30}67 & 21 & 41.2 \\
I-10 Standard Deviation & 17 & 17 & 17 & 16.7 \\
\midrule
\multicolumn{5}{l}{\textit{Level II: Structured Mesh}} \\
II-11 Prefix Sum & \cellcolor{gray!30}84 & \cellcolor{gray!30}80 & 43 & \cellcolor{gray!30}68.8 \\
II-12 Moving Average & 0 & 0 & 0 & 0.0 \\
II-13 Longest Palindrome & 49 & \cellcolor{gray!30}66 & 48 & \cellcolor{gray!30}54.5 \\
II-14 1D Life Game & 27 & 47 & 24 & 32.4 \\
II-15 Pattern Search & 0 & 17 & 17 & 11.1 \\
II-16 Trapping Rain & 33 & 33 & 40 & 35.6 \\
II-17 Diff Array & 48 & \cellcolor{gray!30}62 & 20 & 43.3 \\
II-18 List Ranking & 0 & 0 & 0 & 0.0 \\
II-19 Merge Neighbors & \cellcolor{gray!30}59 & \cellcolor{gray!30}60 & \cellcolor{gray!30}66 & \cellcolor{gray!30}61.8 \\
II-20 Pipeline Hash & \cellcolor{gray!30}52 & 36 & \cellcolor{gray!30}55 & 47.6 \\
\midrule
\multicolumn{5}{l}{\textit{Level III: Global Shuffle}} \\
III-21 Distributed Sort & 33 & 20 & 20 & 24.4 \\
III-22 Median of Medians & 20 & 20 & 0 & 13.3 \\
III-23 Graph Components & 40 & 40 & 14 & 31.3 \\
III-24 BFS Distance & 0 & 0 & 0 & 0.0 \\
III-25 K-Means Iteration & 0 & 0 & 0 & 0.0 \\
III-26 Global Distinct & 33 & 33 & 0 & 22.2 \\
III-27 Collab. Filtering & 0 & 0 & 0 & 0.0 \\
III-28 PageRank Step & 0 & 0 & 0 & 0.0 \\
III-29 Load Balance & 32 & 3 & \cellcolor{gray!30}63 & 32.9 \\
III-30 Matrix Multiply & 0 & 0 & 0 & 0.0 \\
\bottomrule
\end{tabular}
\caption{Success Rate (\%) by task and communication protocol for DeepSeek-V3.1. Tasks with $\geq$50\% success are highlighted in gray background.}
\label{tab:deepseek_detail}
\end{table}


\subsection{Results by Model, Protocol, and Difficulty}

Table~\ref{tab:detailed_all} provides success rates for all model-protocol-difficulty combinations, and Table~\ref{tab:detailed_scale} reports success rates by agent count across models. Together, they confirm that the patterns observed for DeepSeek-V3.1 are consistent across all three models: P2P generally outperforms or matches BP for GPT-OSS-120B and Qwen3, SFS consistently underperforms, and performance degrades monotonically with both complexity level and agent count.

\begin{table}[t]
\centering
\small
\renewcommand{\arraystretch}{1.1}
\begin{tabular}{ll ccc}
\toprule
\textbf{Model} & \textbf{Protocol} & \textbf{L-I} & \textbf{L-II} & \textbf{L-III} \\
\midrule
\multirow{3}{*}{DeepSeek-V3.1}
& BP & 69.7 & 34.5 & 14.7 \\
& P2P & 62.9 & 39.9 & 11.5 \\
& SFS & 53.5 & 30.6 & 9.0 \\
\midrule
\multirow{3}{*}{GPT-OSS-120B}
& BP & 19.6 & 13.9 & 9.7 \\
& P2P & 34.9 & 18.5 & 7.5 \\
& SFS & 27.7 & 11.0 & 9.1 \\
\midrule
\multirow{3}{*}{Qwen3-Next-80B-A3B}
& BP & 9.0 & 3.2 & 1.1 \\
& P2P & 27.6 & 3.4 & 0.3 \\
& SFS & 25.5 & 2.1 & 1.7 \\
\bottomrule
\end{tabular}
\caption{Success Rate (\%) by model, protocol, and difficulty level.}
\label{tab:detailed_all}
\end{table}

\begin{table}[t]
\centering
\footnotesize
\setlength{\tabcolsep}{2pt}
\renewcommand{\arraystretch}{1.0}
\begin{tabularx}{\linewidth}{l *{6}{>{\centering\arraybackslash}X}}
\toprule
\textbf{Model} & \textbf{N=2} & \textbf{N=5} & \textbf{N=10} & \textbf{N=20} & \textbf{N=50} & \textbf{N=100} \\
\midrule
DeepSeek-V3.1 & 61.2 & 48.5 & 39.9 & 33.6 & 19.0 & 18.1 \\
GPT-OSS-120B & 34.4 & 28.0 & 14.0 & 13.2 & 5.2 & 6.4 \\
Qwen3-Next-80B-A3B & 17.2 & 9.1 & 8.6 & 7.4 & 5.1 & 1.3 \\
\bottomrule
\end{tabularx}
\caption{Success Rate (\%) by model and agent count.}
\label{tab:detailed_scale}
\end{table}

\subsection{Communication Density Analysis}

Table~\ref{tab:detailed_density} reports communication density across configurations. A consistent pattern emerges across all models: P2P yields substantially higher densities than BP, reflecting agents' tendency to send multiple targeted messages per pair across rounds; BP densities cluster below 1.0, consistent with one-to-all single broadcasts; and SFS yields notably lower densities than P2P across all models and difficulty levels, indicating that file-based coordination generates sparser cross-agent information flow under our operational definition (read-based transfer counting)---which further explains SFS's systematic under performance on SR despite non-trivial write activity.

\begin{table}[t]
\centering
\small
\renewcommand{\arraystretch}{1.1}
\begin{tabular}{ll ccc}
\toprule
\textbf{Model} & \textbf{Protocol} & \textbf{L-I} & \textbf{L-II} & \textbf{L-III} \\
\midrule
\multirow{3}{*}{DeepSeek-V3.1}
& BP  & 0.56 & 0.82 & 1.46 \\
& P2P & 0.76 & 1.06 & 2.83 \\
& SFS & 0.82 & 1.08 & 1.51 \\
\midrule
\multirow{3}{*}{GPT-OSS-120B}
& BP  & 0.54 & 0.84 & 0.94 \\
& P2P & 1.86 & 2.15 & 2.73 \\
& SFS & 0.73 & 0.54 & 0.95 \\
\midrule
\multirow{3}{*}{Qwen3-Next-80B-A3B}
& BP  & 0.23 & 0.10 & 0.06 \\
& P2P & 1.15 & 0.38 & 0.33 \\
& SFS & 0.48 & 0.07 & 0.09 \\
\bottomrule
\end{tabular}
\caption{Communication Density by model, protocol, and difficulty level.}
\label{tab:detailed_density}
\end{table}


\section{Failure Mode Analysis}
\label{app:failure_analysis}

\subsection{Representative Failure Cases}
\label{app:failure_cases}

Table~\ref{tab:failure_cases} presents representative examples of each failure mode extracted from DeepSeek-V3.1 execution logs. The cases illustrate how failures manifest in practice: premature submission occurs even after reasonable communication volume (Case 4: only 28 of 100 peers contacted before submitting); consensus failure can persist despite near-unanimous agreement, with a single outlier agent preventing full success (Cases 1 and 3); and computation error strikes even when agents have gathered the complete required data (Case 2: off-by-one arithmetic during aggregation).

\begin{table*}[htbp]
\centering
\small
\renewcommand{\arraystretch}{1.2}
\begin{tabularx}{\textwidth}{@{}c >{\raggedright\arraybackslash}p{2.2cm} >{\raggedright\arraybackslash}X >{\raggedright\arraybackslash}X@{}}
\toprule
\textbf{Case} & \textbf{Failure Mode} & \textbf{Description} & \textbf{Key Evidence} \\
\midrule
1 & Consensus Failure & In task I-05, agents communicated extensively but failed to converge, submitting three distinct answers: \{1176, 1182, 1167\}. & 97 of 100 agents submitted 1182; Agent-3 submitted 1176; Agent-55 submitted 1167. \\
\addlinespace
2 & Computation Error & Agent-10 in task I-05 received complete data from all peers but computed an incorrect range count. & Submitted 619 instead of correct answer 620. Arithmetic error during final aggregation. \\
\addlinespace
3 & Consensus Failure & In task I-05 (different instance), 50 agents split between two answers despite communication. & 49 agents submitted 619; Agent-27 submitted 631. \\
\addlinespace
4 & Premature Submission & Agent-77 in task I-06 submitted before collecting sufficient data. & Submitted after receiving data from only 28 of 100 agents. Answer: 208; Expected: 114. \\
\addlinespace
5 & Consensus Failure & 100 agents in task I-06 produced 12 distinct answers for XOR checksum task. & Answers ranged from 42 to 238. Majority (86 agents) converged on 146. \\
\bottomrule
\end{tabularx}
\caption{Representative failure cases from DeepSeek-V3.1 experiments.}
\label{tab:failure_cases}
\end{table*}

\subsection{Failure Mode Definitions and Co-occurrence}

We formally define the three failure modes as follows. \textit{Premature Submission} occurs when an agent invokes \texttt{submit\_result()} before receiving information from a sufficient subset of peers---where ``sufficient'' means the minimum number of agents whose data is required to compute the correct answer. \textit{Consensus Failure} occurs when agents submit multiple distinct answers ($|\{\hat{y}_i\}_{i=1}^{N}| > 1$), indicating that coordination failed to synchronize agents' understanding of global state. \textit{Computation Error} occurs when an agent receives sufficient information but submits an incorrect answer, isolating failures in the reasoning phase from those in the communication phase.

These modes frequently co-occur within single runs, as shown in Table~\ref{tab:failure_cooccurrence}. The high co-occurrence between premature submission and consensus failure (67 cases) suggests a cascading effect: agents submitting early cannot participate in subsequent consensus-building, leaving remaining agents with incomplete information and widening the convergence gap.

\begin{table}[t]
\centering
\small
\renewcommand{\arraystretch}{1.1}
\begin{tabular}{lccc}
\toprule
& \textbf{Premature} & \textbf{Consensus} & \textbf{Compute} \\
\midrule
Premature & 112 & 67 & 45 \\
Consensus & -- & 90 & 52 \\
Compute & -- & -- & 86 \\
\bottomrule
\end{tabular}
\caption{Co-occurrence of failure modes. Diagonal: total occurrences; off-diagonal: joint occurrences.}
\label{tab:failure_cooccurrence}
\end{table}

\subsection{Behavioral Patterns and Leader Emergence}
\label{app:behavior_patterns}

Table~\ref{tab:behavior_patterns_full} compares behavioral metrics between successful and failed runs. Successful runs complete in notably fewer rounds (8.3 vs.\ 12.7), suggesting that effective coordination converges quickly while failed runs engage in extended but ultimately unproductive communication loops. Verification behaviors appear in over 95\% of runs regardless of outcome, confirming that the bottleneck is not communication intent but reasoning quality.

\begin{table}[t]
\centering
\small
\renewcommand{\arraystretch}{1.1}
\begin{tabular}{lcc}
\toprule
\textbf{Metric} & \textbf{Success} & \textbf{Failed} \\
\midrule
Verification Rate & 98.7\% & 95.9\% \\
Strategy Discussion Rate & 93.5\% & 87.2\% \\
Avg.\ Messages per Agent & 31.2 & 27.4 \\
Avg.\ Rounds to Completion & 8.3 & 12.7 \\
\bottomrule
\end{tabular}
\caption{Behavioral comparison between successful and failed runs.}
\label{tab:behavior_patterns_full}
\end{table}

We also examined whether spontaneous leader emergence correlates with task success, classifying an agent as an emergent leader if it receives more than 1.5$\times$ the average number of messages. The results in Table~\ref{tab:leader_emergence_full} are counterintuitive: leader emergence does not consistently improve outcomes, and for Level-III tasks, runs \textit{with} an emergent leader achieve 0\% success versus 33.3\% without one. This suggests that spontaneous centralization at high complexity creates coordination bottlenecks---the designated aggregator becomes overwhelmed by the volume of global data---rather than resolving them.

\begin{table}[t]
\centering
\small
\renewcommand{\arraystretch}{1.1}
\begin{tabular}{lccc}
\toprule
\textbf{Level} & \textbf{Leader Rate} & \textbf{w/ Leader} & \textbf{w/o Leader} \\
\midrule
I & 27.5\% & 56.8\% & 62.1\% \\
II & 21.8\% & 23.5\% & 59.0\% \\
III & 23.8\% & 0.0\% & 33.3\% \\
\bottomrule
\end{tabular}
\caption{Leader emergence rates and associated success rates by complexity level.}
\label{tab:leader_emergence_full}
\end{table}

\subsection{Successful Coordination Examples}
\label{app:success_cases}

To contrast with the failure modes above, we document two illustrative successful patterns. In Case S-1 (Task I-07, $N{=}5$), Agent-0 emerged as coordinator organically: all agents broadcast local results to Agent-0, which computed and rebroadcast the global answer. All agents verified and submitted identically within 4 rounds (100\% success). In Case S-2 (Task I-01, $N{=}10$), agents adopted a distributed verification strategy, with each agent confirming understanding with two neighbors before submission. This redundant verification eliminated consensus failures despite higher message overhead. Both successful patterns share a key property: explicit synchronization checkpoints where agents confirm mutual understanding before proceeding---a discipline entirely absent in failed runs.


\section{Prompt Scaffold Ablation}
\label{app:scaffold_ablation}

To assess the sensitivity of our results to prompt design, we ran a controlled ablation under DeepSeek-V3.1 + P2P with three scaffolding conditions beyond the standard neutral prompt: (a) \textbf{planning round}---a dedicated strategy-discussion round before data exchange begins; (b) \textbf{protocol reminder}---a brief restatement of available communication actions injected at each round; and (c) \textbf{scratchpad hint}---a suggestion to maintain a shared intermediate workspace.

\begin{table}[h]
\centering
\small
\renewcommand{\arraystretch}{1.1}
\begin{tabular}{l ccc}
\toprule
\textbf{Scaffold} & \textbf{L-I SR} & \textbf{L-II SR} & \textbf{L-III SR} \\
\midrule
No scaffold (baseline) & 62.9 & 39.9 & 11.5 \\
+Planning round        & 64.3 & 47.2 & 17.8 \\
+Protocol reminder     & 65.1 & 41.3 & 12.1 \\
+Scratchpad hint       & 63.7 & 44.6 & 14.9 \\
\bottomrule
\end{tabular}
\caption{Prompt scaffold ablation results under DeepSeek-V3.1 + P2P. SR = Success Rate (\%).}
\label{tab:scaffold_ablation}
\end{table}

The planning round yields the most consistent gains ($\sim$5--8\ on Level-II/III); the protocol reminder helps primarily on Level-I; the scratchpad hint benefits intermediate scales ($N{=}10$--20) but cannot prevent collapse at $N{\geq}50$. Critically, qualitative failure patterns remain stable across all conditions: agents continue to communicate actively while failing to translate interaction into correct distributed computation, and the Communication-Reasoning Gap persists regardless of scaffolding. This confirms that the bottleneck reflects genuine LLM limitations in distributed information synthesis rather than a prompting artifact.